% Preface: audience, prerequisites, chapter dependencies, course paths and the code repository.
\paragraph{Who the book is for}
The book is written for two groups of readers. The first is students in Master's programmes in
quantitative finance, financial engineering and investment management who want to see the methods of
their core courses applied to private equity, private credit, real estate, real assets, hedge funds and
structured products. The second is professionals who already work with alternative investments, or are
preparing for the professional examinations that cover them, and who want the mathematics and the code
behind the measures they use every day: the IRR and the public market equivalent, the waterfall, the
loan-to-value and coverage tests, unsmoothed returns, value at risk, and the Sharpe ratio of a short
track record.

\paragraph{What you need to know}
The reader should be comfortable with calculus in several variables, with matrices, inverses and
eigenvalues, and with probability and statistics at the level of a first graduate course: random
variables and their moments, conditional expectation, the normal and lognormal laws, and least squares
regression. No measure theory is assumed; Chapter~3 builds Brownian motion and the It\^o integral from
random walks and states the facts from measure theory that it uses. The code is in Python with NumPy,
SciPy and pandas, and a reader who can run a script and read a short function will be able to follow it.
Nothing in the text depends on R or Octave, which are used only for the independent checks.

\paragraph{How the chapters depend on each other}
The table lists, for each chapter, the earlier chapters whose results it uses. A chapter that is not
listed in a row can be skipped by a reader who wants only that row's chapter.

\begin{center}\small
\begin{tabular}{@{}rlp{0.42\textwidth}@{}}
\toprule
& Chapter & Builds on\\
\midrule
1 & Returns, rates and the shape of risk & none\\
2 & Fund structures, fees and waterfalls & 1\\
3 & Random walks and stochastic calculus & 1\\
4 & Arbitrage, forwards, futures and trees & 1 and 3\\
5 & Black-Scholes & 3 and 4\\
6 & Interest rates, credit and structured products & 4 and 5, with 1 and 3\\
7 & Portfolio theory and factor models & 1\\
8 & Risk and performance measurement & 1 and 7\\
9 & Hedge funds, managed futures and digital assets & 2, 7 and 8, with 5 for the fee as an option\\
10 & Real assets and commercial real estate & 1, 4 and 6\\
11 & Private equity and business valuation & 1 and 2\\
12 & Machine learning in finance & 7 and 8, with 5 for the hedging section\\
\bottomrule
\end{tabular}
\end{center}

\paragraph{Ways through the book}
A course on derivatives and fixed income for alternative assets can take Chapters 1, 3, 4, 5 and 6. A
course on alternative investments can take Chapters 1, 2, 7, 8, 9, 10 and 11, using the statements of
Chapters 4 and 5 where Chapters 9 and 10 call on them. A course on quantitative methods for investment
management can take Chapters 1, 7, 8 and 12. A practitioner who works on private markets may prefer to
start with Chapters 1, 2 and 11 and continue with 10, 8 and 9. Each chapter ends with drills that can be
done by hand, followed by computing projects that use the repository's data and scripts. Worked answers
to the drills follow the exercises, and each project is solved by a function in
\texttt{code/chNN/projects.py} that the tests run, so a reader can check an answer by running the script.

\paragraph{The code repository}
Every number in the text is computed by a script and every figure is drawn by one, so a reader can
change an assumption and see what moves. The repository is public at
\url{https://github.com/ALENZACODES/quant-finance-book}. It has one folder per chapter under
\texttt{code/}, with its scripts, a library module, unit tests, and checks in R and GNU Octave that
recompute the main results by other methods. To set it up, create a Python environment and install the
pinned requirements:
\begin{lstlisting}[language=bash]
python -m venv .venv
.venv/bin/pip install -r requirements.txt
make code CH=06       # numbers, figures and cross-checks for Chapter 6
.venv/bin/python -m pytest -q code/ch06
make book             # the whole book from the committed numbers
\end{lstlisting}
Each script can also be run on its own from the repository root. Random numbers come from fixed seeds
and named streams, so a rerun gives the printed figures exactly. The data snapshots in \texttt{data/}
end in September 2026, and \texttt{data/SOURCES.md} gives the source of each file. A few data sets may
not be redistributed; for those the repository ships a script in \texttt{scripts/} that downloads them
from the original source, and the tests that need them are skipped until they have been fetched. The
computing projects use only the files that ship with the code, read through
\texttt{code/common/shipped.py}, so they run on a fresh copy without any download.

\paragraph{Conventions}
The notation that holds throughout the book, the compounding conventions of each chapter and the
letters that change meaning between chapters are collected under Notation and conventions, just
before Chapter~1.
Results are numbered by chapter, proofs follow the statements, and every simulation names its seed.

From Chapter~6 on, every figure and table carries a tag that says where its numbers come from, and so
do the case studies and the results of the direct lending section of Chapter~6. \kind{exact} marks a closed form or a deterministic calculation from
stated inputs, so a rerun with the same inputs gives the same answer to the last digit.
\kind{simulation} marks a Monte Carlo estimate, which carries sampling error and is reported with its
seed and, where it matters, its standard error. \kind{historical} marks an estimate or a backtest on
the data in the repository, which describes one sample of history and would change with another.
A tag that names two kinds, such as \kind{exact and simulation}, means that the result is shown both
ways, as when a formula is checked against a simulation.
